\GalSourcePageBoundary{075}{L0074}{46}{46}

\begin{center}
{\large\bfseries Application aux équations irréductibles de degré premier.}
\end{center}

\begin{center}{\bfseries PROPOSITION VI.}\end{center}
\GalSourceErrorMark{W09-SE005}

\noindent\textsc{Lemme.} --- {\itshape Une équation irréductible de degré premier ne
peut devenir réductible par l'adjonction d'un radical dont
l'indice serait autre que le degré même de l'équation.}

Car si $r,r',r'',\ldots$ sont les diverses valeurs du radical, et
$Fx=0$ l'équation proposée, il faudrait que $Fx$ se partageât en
facteurs
\[
f(x,r)\times f(x,r')\times\ldots,
\]
tous de même degré, ce qui ne se peut, à moins que $f(x,r)$ ne
soit du premier degré en $x$.

Ainsi, une équation irréductible de degré premier ne peut
devenir réductible, à moins que son groupe ne se réduise à une
seule permutation.

\begin{center}{\bfseries PROPOSITION VII.}\end{center}

\noindent\textsc{Problème.} --- {\itshape Quel est le groupe d'une équation irréductible
d'un degré premier $n$, soluble par radicaux?}

\GalControlReading{D'après la proposition précédente}{W09-CD004}, le plus petit groupe possible,
avant celui qui n'a qu'une seule permutation, contiendra $n$ permutations.
Or, \GalControlReading{un groupe de permutations}{W09-CD005} d'un nombre premier $n$
de lettres ne peut se réduire à $n$ permutations, à moins que
l'une de ces permutations ne se déduise de l'autre par une substitution
circulaire de l'ordre $n$. (\emph{Voir} le Mémoire de M. Cauchy,
\emph{Journal de l'École Polytechnique}, XVII\textsuperscript{e} cahier.) Ainsi, l'avant-dernier
groupe sera
\[
\text{(G)}\qquad
\left\{
\begin{array}{cccccccc}
x_{0},    &x_{1}, &x_{2}, &x_{3}, &\ldots, &x_{n-3}, &x_{n-2}, &x_{n-1}, \\
x_{1},    &x_{2}, &x_{3}, &x_{4}, &\ldots, &x_{n-2}, &x_{n-1}, &x_{0}, \\
x_{2},    &x_{3}, &\ldots, &\ldots, &\ldots, &x_{n-1}, &x_{0}, &x_{1}, \\
\ldots,  &\ldots, &\ldots, &\ldots, &\ldots, &\ldots, &\ldots, &\ldots, \\
x_{n-1}, &x_{0}, &x_{1}, &\ldots, &\ldots, &x_{n-4}, &x_{n-3}, &x_{n-2}.
\end{array}
\right.
\]
$x_{0},x_{1},x_{2},\ldots,x_{n-1}$ étant les racines.
